\chapter{Finite Strings, Induction, and Recursion}
\label{ch:strings}

Formal languages are usually built from finite strings.  This chapter states
the set-theoretic machinery used throughout the reference.

\section{Finite sequences and strings}
\label{sec:strings-01}

\begin{definition}[Finite sequence]
Let \(A\) be a set and let \(n\in\N\).  A \emph{finite sequence in \(A\) of
length \(n\)} is a function \(w:\set{0,\ldots,n-1}\to A\).  Its length is
\(\length(w)=n\).  The unique sequence of length \(0\) is denoted by
\(\epsilon\).
\end{definition}

\begin{definition}[String]
Let \(A\) be an alphabet.  A \emph{string over \(A\)} is a finite sequence in
\(A\).  The set of all strings over \(A\) is
\[
  A^*\defeq\bigcup_{n\in\N}A^n.
\]
\end{definition}

The proof below requires one injection that records two natural numbers in one
natural number.  Fix the metalinguistic abbreviation
\[
  \operatorname{pair}_{\N}(a,b)\defeq
  \frac{(a+b)(a+b+1)}2+b
  \qquad(a,b\in\N).
\]
This map is a bijection \(\N\times\N\to\N\).  Indeed, put
\(t_s=s(s+1)/2\).  For each \(z\in\N\), there is a unique \(s\) with
\(t_s\leq z<t_{s+1}\), because the consecutive finite intervals
\([t_s,t_{s+1})\) partition \(\N\).  Set
\[
  b=z-t_s,
  \qquad
  a=s-b.
\]
The upper inequality gives \(0\leq b\leq s\), so \(a,b\in\N\), and direct
substitution gives \(\operatorname{pair}_{\N}(a,b)=z\).  Conversely, if
\(z=\operatorname{pair}_{\N}(a,b)\), then \(s=a+b\) is the unique interval index and the displayed
inverse recovers \(b\) and then \(a\).  Thus both surjectivity and injectivity
hold.  The abbreviation is used only to encode pairs; it does not add a new
kind of sequence.

